\section{Study Design and Methods}

\subsection{Evidence layers and statistical units}

The A0 registry records dataset, protocol, variant, seed, source and
preprocessing hashes, readout and $K$ provenance, structural source, selection
policy, intervention location, training target, measurement timing, causal
status, artifact availability, label/$K$ isolation, reuse provenance, and
alternative explanations. V1--V22 enter the quantitative atlas. V23 and V24
are registered as boundary evidence but are not pooled with
structural-intervention rows.

The primary retrospective unit is dataset/protocol/readout. Seed and variant
rows are repeated measurements. We do not treat the 1,637 paired records as
1,637 independent data sets. A1 uses raw Delta ARI as its primary descriptive
quantity; headroom-aware summaries and intervention-magnitude descriptors are
sensitivity or post-treatment diagnostics, not causal covariates.

The A1 numbers are imported historical metrics: A1 does not reload labels or
recompute historical ARI values. Consequently, the atlas preserves the
original evaluation and selection provenance, including the historical V21
dataset-selection bias, rather than presenting a newly harmonized label-free
benchmark.

\begin{table}[t]
\centering
\caption{Evidence layers and their inferential boundary.}
\label{tab:layers}
{\small
\setlength{\tabcolsep}{2pt}
\begin{tabular}{@{}p{0.13\linewidth}p{0.25\linewidth}p{0.22\linewidth}p{0.18\linewidth}@{}}
\toprule
Layer & Evidence & Statistical unit & Causal status \\
\midrule
A0/A1 & V1--V22 registry and atlas & dataset / protocol / readout & observational \\
A0 boundary & V23/V24 No-Go records & boundary record & isolated historical boundary \\
E1 & V21 matched N/R/T & dataset; seeds repeated & matched prospective case study \\
E2 & feature and optimizer diagnostics & dataset $\times$ seed & diagnostic/post-hoc \\
E3 & local/global replay boundary & artifact-complete case & observational boundary \\
Phase D & predeclared holdout & six panels & incomplete compute \\
\bottomrule
\end{tabular}
}

\end{table}

\subsection{Failure Atlas and mechanism triage}

A1 admits only artifact-complete cases to offline replay. Metadata-only
snapshots remain provenance records and cannot be reconstructed into replay
evidence. The atlas separates structural opportunity, baseline/headroom,
local/global conversion, and post-treatment magnitude descriptors. Because no
common fixed-graph/null opportunity endpoint exists across all historical
protocols, the atlas is observational rather than a pooled causal test of
structural quality.

A2 constructs a Claim--Evidence Matrix with historical support, counterexamples,
identifiability, required generality, fatal falsifiers, alternative
explanations, novelty, and required new compute. It can retain E1, cancel E1,
or authorize no prospective computation; it cannot invent an E4 after triage.
In the frozen artifacts A2 retained E1 and froze the selection claim, the
primary per-seed endpoint
$S_{\mathrm{full\_ARI}}=\operatorname{ARI}_T-\operatorname{ARI}_R$,
the claim-dependent holdout subset, and the practical threshold $\delta=0.03$.

\subsection{Matched V21 case study}

The prospective case study compares three arms:

\begin{align}
N&=\text{matched None},\\
R&=\text{matched Random assignment policy},\\
T&=\text{topology-dependent selection policy}.
\end{align}

The estimands are
\begin{align}
I_d &= Q(R)-Q(N),\\
S_d &= Q(T)-Q(R),\\
Q(T)-Q(N)&=I_d+S_d.
\end{align}

For E1, $Q$ is the clean-embedding, known-$K$ KMeans ARI for one seed. Thus
$S_{\mathrm{full\_ARI}}$ is the primary seed-level endpoint and
$S_d=\operatorname{mean}_{s\in\{42,123,7\}}S_{\mathrm{full\_ARI},d,s}$ is only
the dataset-level summary used for the four-state rule and figures. Pilot and
confirmation panels are audited and reported separately; they are not pooled.

The losses are frozen as
\begin{align}
L_N &= L_{\mathrm{scMAE}}+\lambda_iL_{\mathrm{InfoMax}},\\
L_R &= L_{\mathrm{scMAE}}+\lambda_iL_{\mathrm{InfoMax}}+\lambda_aL_{\mathrm{JS}}^R,\\
L_T &= L_{\mathrm{scMAE}}+\lambda_iL_{\mathrm{InfoMax}}+\lambda_aL_{\mathrm{JS}}^T.
\end{align}

All arms start from one branchpoint after warmup and Student-$t$ head
initialization and before the first assignment update. Base reconstruction
masks and donors, assignment donors, eligible coordinates, exact effective
budgets, batch order, topology statistics, Gumbel selection noise, and readout
randomness are shared or hash-audited. The Random arm uses zero topology logits
plus the same frozen Gumbel tensor; the Topology arm adds topology-dependent
logits. None preserves the reconstruction and InfoMax terms and optimizer state
but executes no assignment forward and no JS loss.

The primary readout is clean embedding plus known-$K$ KMeans. The Student-$t$
head is a secondary diagnostic. The full label vector is excluded from
preprocessing, graph construction, Gate, loss, optimizer updates, and all model
state transitions. Labels enter only after fitting for outer ARI/NMI evaluation;
benchmark-known $K$ may nevertheless be derived from the outer labels to size
the trainable head and primary readout. E1 is therefore a real-ground-truth,
known-$K$ benchmark, not fully label-free fitting.

The pilot contains cnae9, Mouse Retina, and SMS with seeds $[42,123,7]$; the
confirmation panel contains Baron Human, Campbell, and \texttt{hate\_speech}
with the same seeds. A dataset-level state is Positive or Negative only when
the mean effect exceeds $\pm0.03$ and at least two of three seed effects share
its sign. Observed-Small requires the mean and all three seed effects to be
smaller in absolute value; otherwise the state is Inconclusive. No three-seed
bootstrap is used to claim equivalence.

\subsection{Localization diagnostics}

E2-A compares topology-selected and eligible-but-not-selected coordinates, but
aggregates each metric at dataset-by-seed level. Coordinate-level distributions
are descriptive plots only. Fisher separation, mutual information, and
class-support enrichment are post-hoc label-aware descriptors and are not fit
inputs. E2-B records base, assignment, and InfoMax gradient geometry at T0, T1,
and T2. E2-C clones model, head, and Adam state and executes actual N/R/T Adam
one-step updates. These diagnostics do not establish a universal
objective-compatibility law.

E3 applies the artifact-complete replay gate and keeps label-free geometry
separate from post-hoc supervised neighborhood metrics. The frozen replay
summary has zero eligible rows (\texttt{candidate\_rows=0}), so no E3 replay
was run; the V23 local/global rows remain separate boundary evidence. A local
improvement accompanied by non-positive ARI is therefore only a possible
conversion-failure pattern, not a universal local-to-global theorem.

\subsection{Holdout and closure}

Before E1 outcomes, the holdout candidate pool, input adapters, normalization,
feature selection, graph/model inputs, source hashes, measurement schema, and
claim-dependent activation rule were frozen. The manifest records a scRNA
shortfall and does not silently fill it using outcomes. Phase C selected the
selection claim, so Phase D activated the N/R/T subset and primary
$S_{\mathrm{full\_ARI}}$ endpoint. The expected holdout contained six panels.
No panel produced an evaluable primary endpoint under the frozen dense decoder/
Adam resource path; this is incomplete compute, not an independent replication
result.
